% ============================================================
%  Section 5: OLAP & Business Intelligence Tools
% ============================================================
\section{OLAP \& Business Intelligence}

% -------- subsection: OLAP Definition --------

\begin{frame}
  \frametitle{What is OLAP?}

  \begin{block}{Online Analytical Processing (OLAP)}
    OLAP is a technology (hardware and software) for:
    \begin{itemize}
      \item \textbf{Simple and fast} data analysis
      \item From \textbf{arbitrary perspectives} and at different levels of detail
      \item Without having to define perspectives and levels in advance
    \end{itemize}
    $\Rightarrow$ \textbf{Dynamic analysis of multidimensional data!}
    \hfill\small(based on Codd's 12 OLAP rules, 1993)
  \end{block}

  \begin{block}{The Foundation of OLAP}
    The basis of OLAP is always a \textbf{dimensional data model} ---
    a star or snowflake schema (or a multidimensional database structure
    built from it).
  \end{block}

  \begin{block}{Typical OLAP Questions}
    \begin{itemize}
      \item \textit{What were total sales per region per quarter in 2024?}
      \item \textit{Which product category grew the fastest over the last 3 years?}
      \item \textit{Compare revenue in Germany vs.\ France for Q2 this year.}
    \end{itemize}
  \end{block}
\end{frame}

% ------------------------------------------------------------
\begin{frame}
  \frametitle{Codd's 12 OLAP Rules}

  Edgar F.\ Codd (inventor of the relational model) defined 12 rules
  for a true OLAP system in 1993:

  \begin{columns}[T]
    \column{0.48\textwidth}
    \begin{enumerate}
      \item Multidimensional view of data
      \item Transparent data access
      \item Unrestricted data availability
      \item Consistent performance
      \item Client/server architecture
      \item Support for default dimensions (e.g.\ Time)
    \end{enumerate}

    \column{0.48\textwidth}
    \begin{enumerate}
      \setcounter{enumi}{6}
      \item Dynamic sparsity handling
      \item Multi-user support
      \item Unrestricted operations across all dimensions
      \item Intuitive data manipulation
      \item Flexible reporting
      \item Unlimited dimensions and aggregation levels
    \end{enumerate}
  \end{columns}
\end{frame}

% -------- subsection: OLAP Cube --------

\begin{frame}
  \frametitle{The OLAP Cube}

  \begin{block}{What is an OLAP Cube?}
    An OLAP cube is a \textbf{specialised storage format} with the following properties:
    \begin{itemize}
      \item Matrix-like format optimised for multidimensional access
      \item Enables rapid extraction of pre-summarised data at multiple levels
      \item Organises data in \textbf{hierarchies}
      \item Structures data according to \textbf{dimensions and measures}
    \end{itemize}
  \end{block}

  % PICTURE PLACEHOLDER
  \begin{alertblock}{Picture: 3D OLAP Cube with Axes}
    \textit{Add a 3D cube diagram with three labelled axes:
    Axis 1 = Geography (Country/Region/City),
    Axis 2 = Product Line (Category/Group/Product),
    Axis 3 = Time (Year/Quarter/Month).
    Highlight one cell and label it with an example value:
    ``Geography=Germany, Product=Electronics, Time=2024 Q1: Revenue = \euro{}1.2M''.}
  \end{alertblock}
\end{frame}

% ------------------------------------------------------------
\begin{frame}
  \frametitle{OLAP Cube: From Table to Cube}

  The same data can be represented as a \textbf{relational table} or
  as a \textbf{multidimensional cube} --- they contain the same information.

  % PICTURE PLACEHOLDER
  \begin{alertblock}{Picture: Cube-Table Duality}
    \textit{Add the ``Dualismus Würfel-Tabelle'' diagram:
    Left side: a 3D cube with axes Product, Time, Geography.
    Right side: a flat relational table with columns Article, Branch, Date, Sales.
    Show that the cells in the cube correspond to rows in the table.
    Example rows: Radio/Heidelberg/3.1.2003/7, Radio/Mannheim/3.1.2003/4.}
  \end{alertblock}

  \begin{block}{Tabular Representation}
    OLAP results can be displayed as:
    \begin{itemize}
      \item Simple 2D cross-tables (Rows = Products, Columns = Regions)
      \item Nested tables (multiple dimensions on each axis)
    \end{itemize}
  \end{block}
\end{frame}

% -------- subsection: OLAP Terminology --------

\begin{frame}
  \frametitle{OLAP Terminology}

  \begin{tabularx}{\linewidth}{@{}l X@{}}
    \toprule
    \textbf{Term}       & \textbf{Meaning} \\
    \midrule
    Dimension Name      & The label of a dimension (e.g.\ ``Time'') \\
    Level               & A level of detail within a dimension (e.g.\ Year, Quarter, Month) \\
    Hierarchy           & The arrangement of levels within a dimension \\
    Member              & A single category value within a level (e.g.\ ``Q1 2024'') \\
    Measure             & The numeric value being analysed (e.g.\ Revenue, Quantity) \\
    Crossing            & A specific combination of dimension levels (defines a cell) \\
    N-Way Crossing      & The combination of all lowest levels --- the most granular data \\
    Sparsity            & The proportion of cells in the cube that have no value \\
    Drill-Through       & Access to the underlying transaction-level detail data \\
    \bottomrule
  \end{tabularx}
\end{frame}

% ------------------------------------------------------------
\begin{frame}
  \frametitle{OLAP Terminology: The Time Dimension}

  % PICTURE PLACEHOLDER
  \begin{alertblock}{Picture: Time Dimension with Levels and Hierarchies}
    \textit{Add a dimension tree for the Time dimension:
    Top: Dimension name ``YQ Zeit''.
    Below: two hierarchies side by side.
    Hierarchy 1 (YQ): 2003/2004/2005 $\to$ Q3/Q4/Q1/Q2 (Quarters).
    Hierarchy 2 (YM): 2003/2004/2005 $\to$ Jan/Feb/Mar/Apr/May/Jun/Jul/Aug...
    Label: ``Member'' for individual values (e.g.\ Q3, Jan),
    ``Level'' for the row type, ``Hierarchy'' for the overall branch.}
  \end{alertblock}

  \begin{block}{Key Insight}
    A single dimension can contain \textbf{multiple hierarchies}.
    Example for Time:
    \begin{itemize}
      \item Hierarchy 1 (Calendar): Year $\to$ Quarter $\to$ Month $\to$ Day
      \item Hierarchy 2 (Fiscal): Fiscal Year $\to$ 4-4-5 Month $\to$ Week $\to$ Day
    \end{itemize}
  \end{block}
\end{frame}

% -------- subsection: OLAP Operations --------

\begin{frame}
  \frametitle{OLAP Operation: Roll-Up and Drill-Down}

  \begin{columns}[T]
    \column{0.48\textwidth}
    \begin{block}{Roll-Up (Aggregation)}
      Move \textbf{up} in the dimension hierarchy.\\
      Less detail, more aggregation.\\
      Example:\\
      Day $\to$ Month $\to$ Quarter $\to$ Year
    \end{block}

    \column{0.48\textwidth}
    \begin{block}{Drill-Down (Disaggregation)}
      Move \textbf{down} in the dimension hierarchy.\\
      More detail, less aggregation.\\
      Example:\\
      Year $\to$ Quarter $\to$ Month $\to$ Day
    \end{block}
  \end{columns}

  \vspace{0.5em}
  % PICTURE PLACEHOLDER
  \begin{alertblock}{Picture: Roll-Up / Drill-Down on Product Hierarchy}
    \textit{Add a vertical arrow diagram on the Product hierarchy:
    (Top) Total $\to$ Industry $\to$ Product Group $\to$ Product Family $\to$ (Bottom) Product.
    Left-side arrow pointing up: ``Roll-Up''.
    Right-side arrow pointing down: ``Drill-Down''.}
  \end{alertblock}
\end{frame}

% ------------------------------------------------------------
\begin{frame}
  \frametitle{OLAP Operations: Slice, Dice, and Pivot}

  \begin{block}{Slice --- Fix One Dimension}
    Select a \textbf{single value} on one dimension, reducing the cube by one axis.\\
    \textit{Example:} ``Show me only Q1 2024'' (fixes the Time axis).\\
    Result: a 2D table from a 3D cube.
  \end{block}

  \begin{block}{Dice --- Filter Multiple Dimensions}
    Apply \textbf{multiple filters} across several dimensions, producing a sub-cube.\\
    \textit{Example:} ``Q1 and Q2, Electronics and Clothing, Germany only.''\\
    Result: a smaller cube with restricted ranges.
  \end{block}

  \begin{block}{Pivot (Rotate) --- Swap Axes}
    \textbf{Re-orient} the result by swapping row and column dimensions.\\
    \textit{Example:} Switch from Rows=Products/Columns=Regions
    to Rows=Regions/Columns=Products.\\
    No data changes --- only the presentation perspective changes.
  \end{block}
\end{frame}

% -------- subsection: MOLAP / ROLAP / HOLAP --------

\begin{frame}
  \frametitle{OLAP Storage Variants: Overview}

  \begin{block}{Three Technical Approaches}
    How the multidimensional data is \textbf{physically stored} determines
    performance, scalability, and flexibility.
  \end{block}

  \begin{tabularx}{\linewidth}{@{}l X@{}}
    \toprule
    \textbf{Variant} & \textbf{Storage Type} \\
    \midrule
    \textbf{MOLAP} & Multidimensional (proprietary cube database) \\
    \textbf{ROLAP} & Relational (star/snowflake schema in RDBMS) \\
    \textbf{HOLAP} & Hybrid (combination of both) \\
    \bottomrule
  \end{tabularx}

  % PICTURE PLACEHOLDER
  \begin{alertblock}{Picture: OLAP Scalability Matrix}
    \textit{Add the scalability diagram:
    X-axis = Size of data / tables, Y-axis = Query complexity.
    MDDB OLAP (MOLAP) in bottom-left (small data, complex queries),
    Hybrid OLAP (HOLAP) in middle,
    Relational OLAP (ROLAP) in top-right (large data, simpler queries).
    Use overlapping circles or regions.}
  \end{alertblock}
\end{frame}

% ------------------------------------------------------------
\begin{frame}
  \frametitle{MOLAP: Multidimensional OLAP}

  \begin{block}{How it Works}
    Data is stored in a \textbf{proprietary multidimensional database} (cube file).
    The cube is pre-built and pre-aggregated.
    No traditional table structure --- dimensions and values are stored together.
    All aggregation levels may be pre-computed in one cube.
  \end{block}

  \begin{columns}[T]
    \column{0.48\textwidth}
    \begin{block}{Advantages}
      \begin{itemize}
        \item Optimised for access via dimensions
        \item Very fast for complex queries
        \item Good for many dimensions
        \item Good for deep hierarchies
        \item Good for sparse data (low cardinality)
      \end{itemize}
    \end{block}

    \column{0.48\textwidth}
    \begin{alertblock}{Disadvantages}
      \begin{itemize}
        \item Poor scalability
        \item Problematic for dimensions with very many values (high cardinality)
        \item Size and performance issues at scale
      \end{itemize}
    \end{alertblock}
  \end{columns}
\end{frame}

% ------------------------------------------------------------
\begin{frame}
  \frametitle{ROLAP: Relational OLAP}

  \begin{block}{How it Works}
    Data is stored in a standard \textbf{relational database (RDBMS)}.
    The star schema is used (or fully denormalised tables).
    OLAP cube operations are translated into SQL queries at runtime.
    Views are used to represent different perspectives.
  \end{block}

  \begin{columns}[T]
    \column{0.48\textwidth}
    \begin{block}{Advantages}
      \begin{itemize}
        \item Excellent scalability
        \item Works well for very large datasets (high cardinality)
        \item Flat hierarchies are handled well
        \item Highly dynamic (views can change without rebuilding)
        \item Good for simple queries
      \end{itemize}
    \end{block}

    \column{0.48\textwidth}
    \begin{alertblock}{Disadvantages}
      \begin{itemize}
        \item Slower for very complex queries
        \item Not ideal for many dimensions simultaneously
        \item Deep hierarchies cause many joins
        \item Sparse data is inefficient
      \end{itemize}
    \end{alertblock}
  \end{columns}
\end{frame}

% ------------------------------------------------------------
\begin{frame}
  \frametitle{HOLAP: Hybrid OLAP}

  \begin{block}{The Best of Both Worlds}
    HOLAP combines \textbf{multidimensional} and \textbf{relational} storage.
    The idea is to use each where it performs best:
    \begin{itemize}
      \item \textbf{MOLAP} storage for highly aggregated data and complex queries
      \item \textbf{ROLAP} storage for detailed transaction-level data and large datasets
    \end{itemize}
  \end{block}

  \begin{block}{When to Use HOLAP}
    \begin{itemize}
      \item When you need \textbf{both} summary speed and detail scalability
      \item Complex analytical queries on summaries $\to$ use MOLAP layer
      \item Mass queries on millions of detail rows $\to$ use ROLAP layer
    \end{itemize}
  \end{block}

  \begin{examples}{Practical Assessment}
    HOLAP is \textbf{pragmatic and effective} for most enterprise DWH scenarios.
    The choice of storage type should be driven by the actual query patterns
    and data characteristics --- not by vendor preference.
  \end{examples}
\end{frame}

% -------- subsection: SQL Analytics --------

\begin{frame}[fragile]
  \frametitle{SQL Extensions for OLAP: ROLLUP}

  Standard SQL has powerful OLAP extensions built into most modern databases.

  \begin{block}{\texttt{GROUP BY ROLLUP}}
    Generates subtotals and a grand total along a dimension hierarchy
    in a single query.
  \end{block}

  \begin{lstlisting}[style=sqlstyle]
SELECT region, product, SUM(revenue) AS total_revenue
FROM fact_sales
GROUP BY ROLLUP(region, product);
  \end{lstlisting}

  \begin{block}{Result Structure}
    \begin{itemize}
      \item One row per (region, product) combination
      \item One subtotal row per region (product = NULL)
      \item One grand total row (region = NULL, product = NULL)
    \end{itemize}
    This replaces multiple \texttt{UNION ALL} subqueries with a single,
    optimised pass over the data.
  \end{block}
\end{frame}

% ------------------------------------------------------------
\begin{frame}[fragile]
  \frametitle{SQL Extensions for OLAP: CUBE and GROUPING SETS}

  \begin{block}{\texttt{GROUP BY CUBE}}
    Generates \textbf{all possible combinations} of grouping attributes.
  \end{block}

  \begin{lstlisting}[style=sqlstyle]
SELECT region, product, year, SUM(revenue)
FROM fact_sales
GROUP BY CUBE(region, product, year);
  \end{lstlisting}

  \begin{block}{\texttt{GROUPING SETS}}
    Specify exactly \textbf{which grouping combinations} to compute.
    More flexible than ROLLUP or CUBE.
  \end{block}

  \begin{lstlisting}[style=sqlstyle]
SELECT region, product, year, SUM(revenue)
FROM fact_sales
GROUP BY GROUPING SETS (
    (region, year),   -- by region and year
    (product, year),  -- by product and year
    (region),         -- by region only
    ()                -- grand total
);
  \end{lstlisting}
\end{frame}

% ------------------------------------------------------------
\begin{frame}[fragile]
  \frametitle{SQL Window Functions}

  \begin{block}{What are Window Functions?}
    Window functions perform calculations \textbf{across a set of related rows}
    without collapsing them into a single output row.
    They use the \texttt{OVER()} clause.
  \end{block}

  \begin{lstlisting}[style=sqlstyle]
-- Running total of revenue per region
SELECT
    date_key, region, revenue,
    SUM(revenue) OVER (
        PARTITION BY region
        ORDER BY date_key
        ROWS BETWEEN UNBOUNDED PRECEDING AND CURRENT ROW
    ) AS running_total,
    LAG(revenue) OVER (
        PARTITION BY region ORDER BY date_key
    ) AS prev_period_revenue
FROM fact_sales;
  \end{lstlisting}

  \begin{block}{Common Window Functions}
    \texttt{ROW\_NUMBER()}, \texttt{RANK()}, \texttt{DENSE\_RANK()},
    \texttt{LAG()}, \texttt{LEAD()}, \texttt{NTILE()},
    \texttt{PERCENT\_RANK()}, \texttt{FIRST\_VALUE()}, \texttt{LAST\_VALUE()}
  \end{block}
\end{frame}

% -------- subsection: BI Tools and Reporting --------

\begin{frame}
  \frametitle{From Data to Decision: The BI Layer}

  \begin{block}{What the BI Layer Does}
    The BI layer sits \textbf{on top of} the DWH and OLAP infrastructure.
    It presents data to business users in a consumable format.
    Target users are non-technical: managers, analysts, executives.
  \end{block}

  % PICTURE PLACEHOLDER
  \begin{alertblock}{Picture: Complete BI Stack}
    \textit{Add a vertical stack diagram from bottom to top:
    (1) Source Systems (OLTP, SAP, CRM)
    $\to$ (2) ETL / ELT Pipeline
    $\to$ (3) Data Warehouse (Staging, Core, Data Marts)
    $\to$ (4) OLAP / Semantic Layer (Cube or SQL views)
    $\to$ (5) BI Tools (Dashboards, Reports, Ad-hoc queries)
    $\to$ (6) Business Decision.
    Add example product names at each layer (e.g.\ Airflow, Snowflake, Power BI).}
  \end{alertblock}
\end{frame}

% ------------------------------------------------------------
\begin{frame}
  \frametitle{Types of BI Analysis}

  \begin{block}{Descriptive Analytics --- \textit{What happened?}}
    Summarises historical data. Standard reports and dashboards.\\
    \textit{Example:} Monthly revenue report.
  \end{block}

  \begin{block}{Diagnostic Analytics --- \textit{Why did it happen?}}
    Drill-down analysis to identify root causes.\\
    \textit{Example:} Revenue dropped in Q3 --- which region? Which product?
  \end{block}

  \begin{block}{Predictive Analytics --- \textit{What will happen?}}
    Statistical models and ML to forecast future outcomes.\\
    \textit{Example:} Demand forecasting, churn prediction.
  \end{block}

  \begin{block}{Prescriptive Analytics --- \textit{What should we do?}}
    Recommends actions based on data and models.\\
    \textit{Example:} Automated pricing engine, inventory optimisation.
  \end{block}
\end{frame}

% ------------------------------------------------------------
\begin{frame}
  \frametitle{Types of BI Reports and Dashboards}

  \begin{block}{Standard / Operational Reporting}
    \begin{itemize}
      \item Fixed-format, scheduled reports delivered regularly
      \item Example: End-of-day sales summary, monthly P\&L
      \item Designed once, run repeatedly
    \end{itemize}
  \end{block}

  \begin{block}{Ad-Hoc Analysis}
    \begin{itemize}
      \item Users create their own queries and reports on demand
      \item Requires a user-friendly \textbf{semantic layer}
      \item No IT involvement needed for each analysis
    \end{itemize}
  \end{block}

  \begin{block}{Interactive Dashboards}
    \begin{itemize}
      \item Visual, real-time display of KPIs
      \item Users can filter, drill down, and explore data
      \item Tools: \textbf{Power BI}, \textbf{Tableau}, \textbf{Looker},
        \textbf{Metabase}, \textbf{SAP BusinessObjects}
    \end{itemize}
  \end{block}
\end{frame}

% ------------------------------------------------------------
\begin{frame}
  \frametitle{KPIs and Metrics}

  \begin{block}{Key Performance Indicator (KPI)}
    A quantifiable measure used to evaluate success in achieving
    a business objective. Not all metrics are KPIs --- only those
    that are directly tied to strategic goals.
  \end{block}

  \begin{block}{Properties of a Good KPI}
    \begin{itemize}
      \item \textbf{Specific}: tied to a concrete business goal
      \item \textbf{Measurable}: expressed as a number
      \item \textbf{Actionable}: something the business can actually influence
      \item \textbf{Timely}: available when decisions need to be made
      \item \textbf{Consistently defined}: same meaning across all departments
    \end{itemize}
  \end{block}

  \begin{examples}{Example KPIs}
    Monthly Recurring Revenue (MRR), Customer Acquisition Cost (CAC),
    Net Promoter Score (NPS), Inventory Turnover Rate,
    Customer Lifetime Value (CLV), Sales Conversion Rate.
  \end{examples}
\end{frame}

% ------------------------------------------------------------
\begin{frame}
  \frametitle{The Semantic Layer}

  \begin{block}{Definition}
    The \textbf{semantic layer} (business layer) maps the physical DWH
    tables and columns to \textbf{business-friendly terms} that
    non-technical users can understand and trust.
  \end{block}

  \begin{block}{What the Semantic Layer Provides}
    \begin{itemize}
      \item Abstracts SQL complexity from business users
      \item Defines \textbf{metrics, KPIs, and dimensions} in business language
      \item Ensures consistent definitions company-wide:
        ``Revenue'' always means the same thing everywhere
      \item Enables self-service analytics without IT involvement
    \end{itemize}
  \end{block}

  \begin{examples}{Example Products}
    \textbf{SSAS Tabular} (Microsoft), \textbf{Looker LookML},
    \textbf{dbt metrics layer}, \textbf{Power BI datasets},
    \textbf{SAP BW BEx} (Business Explorer)
  \end{examples}
\end{frame}

% ------------------------------------------------------------
\begin{frame}
  \frametitle{Summary: From Raw Data to Business Decision}

  % PICTURE PLACEHOLDER
  \begin{alertblock}{Picture: End-to-End Flow Diagram}
    \textit{Add a horizontal arrow diagram showing the complete journey:
    (1) Raw Transactions (OLTP/ERP)
    $\to$ (2) ETL/ELT with Data Quality
    $\to$ (3) Data Warehouse (Staging + Core + Data Mart)
    $\to$ (4) OLAP Cube / Semantic Layer
    $\to$ (5) BI Dashboard / Report
    $\to$ (6) Business Decision.
    Add small icons or example artefacts at each step.}
  \end{alertblock}

  \begin{block}{Key Takeaways}
    \begin{itemize}
      \item \textbf{BI} is the ``intelligence service'' of the enterprise
      \item The \textbf{DWH} is the foundation: integrated, historical, non-volatile
      \item \textbf{ETL and data quality} determine the value of all analytics
      \item \textbf{OLAP} enables fast, flexible, multidimensional analysis
      \item \textbf{BI tools} translate data into decisions for business users
    \end{itemize}
  \end{block}
\end{frame}
